\documentclass[12pt]{article}

\usepackage[portuguese]{babel}
\usepackage[utf8]{inputenc}
\usepackage[T1]{fontenc}
\usepackage{geometry}
\usepackage{array}
\usepackage{tabularx}
\usepackage{colortbl}
\usepackage{xcolor}
\usepackage{fancyhdr}
\usepackage{lastpage}

\geometry{margin=2.5cm}

% Cor cinzenta para cabeçalhos de tabela
\definecolor{tablegray}{gray}{0.9}
\definecolor{primaryblue}{RGB}{0, 51, 102}

% Cabeçalho e rodapé
\pagestyle{fancy}
\fancyhf{}

\fancyhead[L]{Registo de Reunião | EcoRide Solutions}
\fancyhead[R]{\thepage\ de \pageref{LastPage}}

\renewcommand{\headrulewidth}{0.4pt}
\renewcommand{\footrulewidth}{0pt}

\begin{document}

{\noindent\LARGE\bfseries\color{primaryblue} Ata de Reunião | Requisitos não-funcionais}

\vspace{0.4cm}

\section*{Identificação da Reunião}

\begin{tabularx}{\textwidth}{|l|X|}
\hline
\rowcolor{tablegray}
\textbf{ID Reunião} & \#007 \\ \hline
\textbf{Tipo de reunião} & Levantamento de requisitos não-funcionais\\ \hline
\textbf{Data} & 13 de março de 2026 \\ \hline
\textbf{Hora Início} & 9:00 \\ \hline
\textbf{Hora Fim} & 11:00 \\ \hline
\textbf{Local/Meio} & Chamada de videoconferência \\ \hline
\textbf{Estado} & Finalizada \\ \hline
\end{tabularx}

\vspace{0.5cm}

\section*{Participantes}

\begin{tabularx}{\textwidth}{|l|l|X|}
\hline
\rowcolor{tablegray}
\textbf{Nome} & \textbf{Função} & \textbf{Entidade} \\ \hline
Duarte Escairo & Programador & Equipa de Desenvolvimento \\ \hline
Eduardo Fernandes & Analista & Equipa de Desenvolvimento \\ \hline
Inês Ferreira & Coordenadora & Equipa de Desenvolvimento \\ \hline
Luís Soares & Analista & Equipa de Desenvolvimento \\ \hline
Tiago Figueiredo & Programador & Equipa de Desenvolvimento \\ \hline
\end{tabularx}

\vspace{0.8cm}
\hrule % Linha separadora opcional
	\vspace{0.2cm}

\section*{Objetivo da Reunião}

A reunião teve como objetivo identificar, discutir e formalizar os requisitos não‑funcionais do sistema de gestão da EcoRide Solutions. Pretendeu‑se garantir que, para além das funcionalidades já levantadas, o sistema cumpra critérios de desempenho, segurança, usabilidade, fiabilidade, portabilidade e manutenção adequados ao contexto operacional da oficina.

Durante a sessão, foram analisadas práticas de empresas semelhantes, bem como limitações observadas nos processos atuais, de forma a assegurar que o sistema final apresenta níveis de qualidade compatíveis com as necessidades reais dos utilizadores.

\vspace{0.5cm}

\section*{Principais Problemas Identificados}

Durante a reunião, a equipa analisou em detalhe os aspetos não‑funcionais que o sistema deverá cumprir para garantir um funcionamento eficiente, seguro e adequado ao contexto operacional da EcoRide Solutions. Começou‑se por discutir a importância do desempenho, destacando‑se a necessidade de tempos de resposta reduzidos nas operações mais frequentes, como a criação de ordens de serviço, a pesquisa de clientes e a consulta do estado das reparações. Foi igualmente salientado que o sistema deverá suportar vários utilizadores em simultâneo sem degradação perceptível da performance, tendo em conta que secretária, mecânicos e gerente poderão estar a utilizá‑lo ao mesmo tempo.

Abordaram‑se questões de usabilidade e foram discutidas boas práticas de design baseadas nas heurísticas de Nielsen, reconhecendo‑se que o sistema deve ser intuitivo e fácil de aprender, de forma a não comprometer o ritmo de trabalho da oficina.

No que diz respeito à segurança, foi consensual que o sistema deverá implementar perfis de acesso distintos para gerente, secretária e mecânicos, assegurando que cada utilizador apenas acede às funcionalidades que lhe dizem respeito. Foram também discutidas medidas de proteção de dados sensíveis, como encriptação de informação pessoal e registo de todas as ações críticas para efeitos de auditoria.

A equipa discutiu também requisitos relacionados com portabilidade e compatibilidade, sublinhando que o sistema deverá funcionar corretamente nos principais navegadores e em dispositivos de diferentes resoluções, como os tablets utilizados na oficina. Por fim, foram abordados aspetos de manutenibilidade, defendendo‑se uma arquitetura modular que facilite futuras alterações e simplifique a gestão.

\vspace{0.5cm}

\section*{Requisitos obtidos}
	
\begin{tabularx}{\textwidth}{@{}|l|X|@{}}
\hline
\rowcolor{gray!10} \textbf{Número} & \textbf{Descrição} \\ \hline


\textbf{1} & O sistema deve ser modular, permitindo a adição, substituição ou atualização de componentes sem impacto nas restantes funcionalidades. \\ \hline

\textbf{2} & O acesso ao sistema deve ser autenticado com credenciais individuais. \\ \hline

\textbf{2} & Devem existir perfis de acesso distintos, com permissões específicas. \\ \hline

\textbf{4} & Dados sensíveis devem ser armazenados de forma encriptada. \\ \hline

\textbf{5} & O sistema deve impedir a entrega de uma trotinete caso existam pagamentos pendentes. \\ \hline

\textbf{6} & A interface deve ser simples e intuitiva para permitir que novos funcionários aprendam a utilizá‑la rapidamente. \\ \hline

\textbf{7} & Os formulários devem validar automaticamente campos obrigatórios e formatos. \\ \hline

\textbf{8} & O sistema deve minimizar o número de cliques necessários para operações frequentes. \\ \hline

\textbf{9} & O sistema deve apresentar listas de opções e/ou autocompletar campos sempre que possível. \\ \hline

\textbf{10} & Todas as ações críticas devem ser registadas com data, hora e utilizador. \\ \hline

\textbf{11} & O sistema deve garantir que duas pessoas não podem editar a mesma ordem de serviço em simultâneo. \\ \hline

\textbf{12} & O sistema deve funcionar em qualquer navegador moderno e ser compatível com diferentes resoluções. \\ \hline

\textbf{13} & O sistema deve permitir criar uma nova ordem de serviço em menos de 5 segundos, mesmo em períodos de maior carga. \\ \hline

\textbf{14} & As pesquisas por cliente, trotinete ou número de ordem devem apresentar resultados em menos de 2 segundos. \\ \hline

\textbf{15} & O sistema deve suportar, sem degradação significativa, todos os funcionários em simultâneo. \\ \hline

\textbf{16} & Atualizações devem ser planeadas fora do horário de funcionamento da oficina. \\ \hline

\end{tabularx}

\vspace{0.5cm}

\section*{Próximos Passos}

\begin{itemize}
\item Incorporar os requisitos não‑funcionais no documento final de requisitos.
\item Rever a estrutura geral dos requisitos para garantir consistência entre funcionais e não‑funcionais.
\item Validar o documento consolidado com todos os stakeholders.
\item Utilizar estes requisitos como base para a modelação da arquitetura e definição dos testes de qualidade.
\end{itemize}

\vspace{0.5cm}

\section*{Encerramento}

A reunião foi encerrada após validação dos requisitos não‑funcionais levantados e definição dos próximos passos. A equipa de desenvolvimento comprometeu‑se a integrar estes requisitos na documentação oficial e a garantir que serão considerados nas fases seguintes do projeto.

\end{document}